\begin{proofbox}
	\textbf{\textcolor{CProof}{思路提示}}

	要证明$a\parallel b$，只需证明$a$与$b$在同一平面$\gamma$内且没有公共点，然后由平面内两直线无公共点则平行得出结论。

	\medskip
	\textbf{\textcolor{CProof}{正式推导}}
	\begin{description}[style=nextline, leftmargin=2.8em, labelsep=0.5em, itemsep=0.55em]
		\item[\textbf{步骤一（共面判定）：}] 因为$\gamma\cap\alpha=a$，$\gamma\cap\beta=b$，所以$a\subset\gamma$，$b\subset\gamma$。
			\par\textbf{理由：} 两个平面相交，交线属于这两个平面，因此$a,b$都在$\gamma$内。

		\item[\textbf{步骤二（无公共点证明）：}] 由于$\alpha\parallel\beta$，故$\alpha$与$\beta$没有公共点。若$a$与$b$相交于某点$P$，则$P\in a\subset\alpha$且$P\in b\subset\beta$，$P$将成为$\alpha,\beta$的公共点，矛盾。因此$a$与$b$无公共点。
			\par\textbf{理由：} 面面平行定义：两平面无公共点。

		\item[\textbf{步骤三（平行结论）：}] 在平面$\gamma$内，两条直线$a,b$没有公共点，所以$a\parallel b$。
			\par\textbf{理由：} 同一平面内，两条直线不相交即平行。
	\end{description}

	\medskip
	\textbf{\textcolor{CProof}{结论回扣}}

	因此，若$\alpha\parallel\beta$，且$\gamma$与$\alpha,\beta$分别交于$a,b$，则$a\parallel b$。这就是面面平行性质定理。
\end{proofbox}
